\chapter{Preliminaries}\label{preliminaries}
This \emph{optional} chapter contains the stuff that your reader needs to know in order to understand
your work.
Your ``audience'' consists of fellow third year computing science bachelor students who have done the same
core courses as you have, but not necessarily the same specialization, minor, or free electives.

You do not have to call this chapter \emph{Preliminaries} and it does not need to be limited to one chapter.
Some preliminary topics deserve their own chapter.

\textbf{Preliminaries:}
\begin{itemize}
    \item LLM Hallucination
    \item Cross-Lingual Natural Language Inference (NLI)
    \item Open Web Index (OWI) Architecture \todo[color=gray,nolist]{If i decide on using it.}
\end{itemize}


\section{Hallucination in Large Language Models}\label{pre:sec:llm}
In this thesis, we will be looking at the output of different Large Language Models (LLM). LLMs are massive artificial neural networks
pre-trained on vast amounts of text data, designed to understand and generate human language. They show strong understanding in natural language and solving 
complex tasks.  \cite{zhao2023survey} While these outputs often contain useful and correct information, they are prone to \textit{hallucinating}. 
Hallucination of a LLM consists of the model generating content that is not faithful to its source. There exist multiple types of hallucinating, 
and following the taxonomy proposed by
\cite{huang2025survey}, we can categorise these into two types: 
\begin{itemize}
    \item \textbf{Factuality Hallucinations} 
    \\ Discrepancies between generated content and verifiable real-world facts. 
    This includes \textit{Fabrications} (made-up claims) and \textit{Contradictions} (statements that directly oppose reality).
    \item \textbf{Faithfulness Hallucinations} 
    \\ Divergence from the provided user input or context. This includes \textit{Instruction} (following the wrong instructions from the prompt), 
    \textit{Context} (Not understanding the context) and \textit{Logical} (Incorrect logic, e.g. math errors).  
    In information extraction (IE) tasks, these often manifest as ``Missing Spans" or ``Incorrect Types" \cite{gao2024empirical}.
\end{itemize}

\noindent Classifying the LLM Hallucination, allows us to have insight on what went wrong and where to improve. It gives us the ability to do a system-level
diagnosis of the pipline \cite{vinay2025failuremodesllmsystems}. For instance, a \textit{faithfulness} error indicates a failure to maintain ``context-boundary" stability \cite{vinay2025failuremodesllmsystems}, 
where a \textit{factuality} error might need external grounding through retrieval. Because LLMs are often unable to faithfully detect their own factual errors 
  \cite{chen2023felm}, deterministic verification is required to ensure airtight verification. 



\section{Cross-Lingual Natural Language Inference}\label{pre:sec:nli}
Natural Language Inference (NLI), also known as Recognizing Textual Entailment (RTE), 
is a core task in Natural Language Processing (NLP) \cite{conneau2020unsupervised}. 
NLP is a multidisciplinary field at the crossover between Computing Science, Linguistics, and Artificial Intelligence, where scientists focus on 
helping computers understand natural language. Common applications consist of machine translation, 
sentiment analysis, and automated summarisation. The latter is 
highly relevant for this thesis, since we deal with Information Extraction (IE), specifically
 extracting text and entities from websites \cite{gao2024empirical}, and Semantic Validation, 
 which involves logically validating the semantic meaning of different sentences to detect errors \cite{huang2025survey, chen2023felm}.

The task of NLI in NLP is determining the logical relationship between two pieces of text: 
a \textit{premise} and a \textit{hypothesis} \cite{conneau2020unsupervised}. Based on the truth of the premise, 
the model assigns one of three categorical labels as demonstrated in Table \ref{tab:nli_examples}:

\begin{table}[h]
    \centering
    \caption{Examples of Natural Language Inference (NLI) relationships.}
    \label{tab:nli_examples}
    \begin{tabular}{lll}
        \toprule
        \textbf{Premise} & \textbf{Hypothesis} & \textbf{NLI Label} \\
        \midrule
        Frank Boeijen is from Nijmegen. & Nijmegen is the hometown of Frank Boeijen. & \textit{Entailment} \\
        Frank Boeijen is from Nijmegen. & Frank Boeijen is from Amsterdam. & \textit{Contradiction} \\
        Frank Boeijen is from Nijmegen. & Frank Boeijen is coming back to Nijmegen. & \textit{Neutral} \\
        \bottomrule
    \end{tabular}
\end{table}

\begin{itemize}
    \item \textbf{Entailment:} The hypothesis is guaranteed to be true if the premise is true. In Table \ref{tab:nli_examples}, 
    the premise logically necessitates that Nijmegen is Boeijen's hometown \cite{conneau2020unsupervised}.
    \item \textbf{Contradiction:} The hypothesis is guaranteed to be false if the premise is true. 
    Since the premise establishes Nijmegen as the origin, a claim of Amsterdam constitutes a direct contradiction \cite{conneau2020unsupervised}.
    \item \textbf{Neutral:} The truth of the hypothesis cannot be determined from the premise alone. 
    While Frank Boeijen may be ``coming back" to Nijmegen, the original premise regarding his origin does not provide sufficient
     information to confirm or deny a return trip \cite{conneau2020unsupervised}.
\end{itemize}

\noindent In the context of this thesis, we utilize \textit{Cross-Lingual} NLI to bridge the language gap 
between Dutch source text and English generated summaries. By using a multilingual model like XLM-RoBERTa, 
we can evaluate a Dutch premise against an English hypothesis without an intermediate translation step,
 preserving the semantic nuances of the original concert metadata \cite{conneau2020unsupervised}.